\section{Conclusions}\label{sec10}
Here we present a short guide in this report. In other words, we
are going to answer the following question. What should one do if
one likes to calculate one-loop counterterms, anomalies, an
expansion term in the effective action, or something else which is
defined by the heat kernel expansion? The first step is to find
the bulk part of the variation of the classical action
(\ref{actquad}) and corresponding operator $D$. Next one has to
bring this operator to the canonical form (\ref{laplaceb}) with
the help of (\ref{D1}), (\ref{oab}) and (\ref{Eab}). However, one
can first consult sec.\ \ref{sec3}. Probably, relevant expressions
can be found there. If the problem in question contains neither
boundaries nor singularities, one can look in sec.\ \ref{sec4} for
an expression for the heat kernel coefficient, or for a reference,
or for a method. In the case of the boundaries one has to proceed
with sec.\ \ref{sec5}, in the case of singularities -- with sec.\
\ref{sec6}. Relations between the heat kernel coefficients and
quantum anomalies can be found in sec.\ \ref{sec7}. The results
going beyond the standard heat kernel expansion are collected in
sec.\ \ref{sec8}. In particular, in this section we explain how
one can extract leading non-localities from the effective action,
and what the heat kernel looks like if the background is
approximately covariantly constant (in this context we also consider
invariant operators on homogeneous spaces). Exact results for the
effective action which can be obtained with the help of the heat
kernel expansion are reviewed in sec.\ \ref{sec9}.

The Casimir energy is one of the most ``classical'' applications
of the heat kernel and zeta function technique (cf. 
\cite{Blau:1988kv}). It follows 
from the locality of the heat kernel expansion that the divergences
in the Casimir energy are given by volume and surface integrals
of some local invariants. Therefore, if the boundaries are being
moved in such a way that the boundary values of the background fields 
remain unchanged, the boundary contributions
to the divergences also remain constant. This leads to the well
known conclusion \cite{Bordag:2001qi} that there are no boundary
divergences in the Casimir force (which is defined, roughly speaking
as a variation of the Casimir energy under infinitesimal translations
of the boundary). Consequently, one can assign a unique value to
the Casimir force (see, e.g., \cite{Milton:2002vm}). This observation,
however does not mean that the quantum field theory on a manifold
with boundary is finite. In general, some surface counterterms
are required (at one-loop order they may be read off from 
sec.\ \ref{sec5}). Moreover, if the background field are non-trivial,
the boundary divergences will not be constant. Similar arguments
created certain scepticism towards reliability of the Casimir energy
calculation \cite{Graham:2002xq}. This point has not been settled
so far for a ``generic'' theory. We may add that in supersymmetric 
theories cancellations between divergences in the bosonic and
fermionic sectors appear if the boundary terms are considered
together with the volume terms \cite{Bordag:2002dg}, therefore,
separation of boundary and volume contributions is not always natural
for that theories. Renormalization of self-interacting theories
on manifolds with boundaries was considered in 
\cite{Symanzik:1981wd,Odintsov:1990ny,
McAvity:1993fq,Wiesendanger:1994mw,Barvinsky:1996dp}
where one can find further references. Some aspects of the
relationship between the Casimir energy calculations and the
heat kernel coefficients have been clarified recently by
Fulling \cite{Fulling:2003zx}.

Of course, not everything can be found in this report.
Several topics are very close to the subject of this review,
but are not included.
\begin{enumerate}
\item
The heat kernel expansion can be successfully applied to quantum field
theory at finite temperature \cite{Dowker:1978md}. A new interesting
development in this field is related to the so called non-linear
spectral problem \cite{Fursaev:2002vi} (see \cite{Fursaev:2001yu}
for an overview).
\item
The heat kernel expansion has interesting applications to
integrable models and, in particular, to the
KdV hierarchies (see \cite{Das:1989fn} for an elementary introduction).
\item
Recently some attention has been attracted 
\cite{Dowker:2000bi,Dowker:2000vm,Seeley:2001,Avramidi:2001aj}
to the so called N/D or Zaremba problem which appears when one
defines Neumann and Dirichlet boundary conditions on two
(intersecting) components of the boundary. It is unclear whether
this problem may have applications to quantum theory.
\item
Instead of considering asymptotics of the heat trace ${\rm Tr} (fe^{-tD})$
one can also consider an asymptotic expansion for individual matrix
elements of the heat kernel $\langle f_1, e^{-tD} f_2 \rangle$ (which are
called the heat content asymptotics since they remind short time
asymptotics of the total heat content in a manifold with the specific heat
$f_1$ and the initial temperature distribution $f_2$). Such asymptotics do
not contain negative powers of $t$. More details can be found in
\cite{Kirsten:2001wz}.
\item
Many results on the heat kernel asymptotics can be extended
to higher-order differential operators (see, e.g., 
\cite{Gilkey:1980,Gusynin:1990zt}) and to differential operators
in superspace \cite{Buchbinder:1986im}.
\end{enumerate}

Although quite a lot is already known about the heat kernel expansion,
many interesting problems still remain open. There are many 
opportunities to extend and generalise the results presented in this
report. This refers especially to the material of sec.\ 
\ref{sec5}-\ref{sec9} where one could add new types of the operators,
boundary conditions, geometries, and singularities. 
There is a completely new field of research related to the
heat kernel expansion where very little has been done so far.
This is 
an extension to non-commutative geometry. This  problem
is an especially intriguing one since one can expect very unusual
properties of the spectral functions because of very unusual properties
of corresponding field theories in the ultra violet asymptotics.


\section*{Acknowledgements}
I am grateful to my collaborators S.~Alexandrov, M.~Bordag,
T.~Branson, E.~Elizalde, H.~Falomir, P.~Gilkey, D.~Grumiller,
K.~Kirsten,
W.~Kummer, H.~Liebl, V.~Lyakhovsky, V.~Marachevsky,
Yu.~Novozhilov, M.~Santangelo, N.~Shtykov, P.~van~Nieuwenhuizen
and A.~Zelnikov who contributed in many ways to the material
presented in this report. I have benefited from enlightening discussions
with I.~Avramidi, A.~Barvinsky, G.~Esposito, D.~Fursaev, G.~Grubb,
A.~Kamenshchik, V.~Nesterenko and 
with my friends and colleagues at St.~Petersburg,
Leipzig and Vienna.

I am grateful to all readers who suggested their comments on the
previous version of the manuscript. I also thank the referee
for useful critical remarks.

This work has been supported by Project BO 1112/12-1 of the
Deutsche Forschungsgemeinschaft. I am grateful to E.~Zeidler
for his kind hospitality at the Max Planck Institute for 
Mathematics in the Sciences where a part of this work has
been done.


